%%%%%%%%%%%%%%%%%%%%%%%%%%%%%%%%%%%%%%%%%%
% Short Sectioned Assignment LaTeX Template Version 1.0 (5/5/12)
% This template has been downloaded from: http://www.LaTeXTemplates.com
% Original author:  Frits Wenneker (http://www.howtotex.com)
% License: CC BY-NC-SA 3.0 (http://creativecommons.org/licenses/by-nc-sa/3.0/)
%%%%%%%%%%%%%%%%%%%%%%%%%%%%%%%%%%%%%%%%%

% \documentclass[paper=a4, fontsize=11pt]{scrartcl} % A4 paper and 11pt font size
\documentclass[11pt, a4paper]{book}
\usepackage[T1]{fontenc} % Use 8-bit encoding that has 256 glyphs
\usepackage[utf8]{inputenc}
\usepackage{fourier} % Use the Adobe Utopia font for the document - comment this line to return to the LaTeX default
\usepackage{listings} % para insertar código con formato similar al editor
\renewcommand{\lstlistingname}{Código}
\lstset{
  basicstyle=\ttfamily,
  columns=fullflexible,
  frame=single,
  breaklines=true,
  postbreak=\mbox{\textcolor{red}{$\hookrightarrow$}\space},
}
\usepackage[spanish, es-tabla]{babel} % Selecciona el español para palabras introducidas automáticamente, p.ej. "septiembre" en la fecha y especifica que se use la palabra Tabla en vez de Cuadro
\usepackage{url} % ,href} %para incluir URLs e hipervínculos dentro del texto (aunque hay que instalar href)
\usepackage{graphics,graphicx, float} %para incluir imágenes y colocarlas
\graphicspath{{images/}{../images/}}
\usepackage[gen]{eurosym} %para incluir el símbolo del euro
%\usepackage{cite} %para incluir citas del archivo <nombre>.bib
\usepackage{enumerate}
\usepackage{hyperref}
\usepackage{graphicx}
\usepackage{tabularx}
\usepackage{booktabs}
\usepackage{subfiles}
\usepackage[style=numeric,sorting=none]{biblatex}
\addbibresource{bibliografia.bib}

\usepackage[table,xcdraw]{xcolor}
\hypersetup{
	colorlinks=true,	% false: boxed links; true: colored links
	linkcolor=black,	% color of internal links
	urlcolor=cyan		% color of external links
}
\renewcommand{\familydefault}{\sfdefault}
\usepackage{fancyhdr} % Custom headers and footers
\pagestyle{fancyplain} % Makes all pages in the document conform to the custom headers and footers
\fancyhead[L]{} % Empty left header
\fancyhead[C]{} % Empty center header
\fancyhead[R]{Adrián Álvarez Sáez} % My name
\fancyfoot[L]{} % Empty left footer
\fancyfoot[C]{} % Empty center footer
\fancyfoot[R]{\thepage} % Page numbering for right footer
%\renewcommand{\headrulewidth}{0pt} % Remove header underlines
\renewcommand{\footrulewidth}{0pt} % Remove footer underlines
\setlength{\headheight}{13.6pt} % Customize the height of the header

\usepackage{titlesec, blindtext, color}
\definecolor{gray75}{gray}{0.75}
\newcommand{\hsp}{\hspace{20pt}}
\titleformat{\chapter}[hang]{\Huge\bfseries}{\thechapter\hsp\textcolor{gray75}{|}\hsp}{0pt}{\Huge\bfseries}
\setcounter{secnumdepth}{4}
\setcounter{tocdepth}{3}
\usepackage[Lenny]{fncychap}
\usepackage{csquotes}



\begin{document}

	% Plantilla portada UGR
	\begin{titlepage}
\newlength{\centeroffset}
\setlength{\centeroffset}{-0.5\oddsidemargin}
\addtolength{\centeroffset}{0.5\evensidemargin}
\thispagestyle{empty}

\noindent\hspace*{\centeroffset}\begin{minipage}{\textwidth}

\centering
\includegraphics[width=0.9\textwidth]{logos/logo_ugr.jpg}\\[1.4cm]

\textsc{ \Large TRABAJO FIN DE GRADO\\[0.2cm]}
\textsc{ GRADO EN INGENIERÍA INFORMÁTICA}\\[1cm]

{\Huge\bfseries Título \\}
\noindent\rule[-1ex]{\textwidth}{3pt}\\[3.5ex]
{\large\bfseries Subtítulo }
\end{minipage}

\vspace{2.5cm}
\noindent\hspace*{\centeroffset}
\begin{minipage}{\textwidth}
\centering

\textbf{Autor}\\ {Adrián Álvarez Sáez}\\[2.5ex]
\textbf{Director}\\ {Juan Julián Merelo Guervós}\\[2cm]
\includegraphics[width=0.3\textwidth]{logos/etsiit_logo.png}\\[0.1cm]
\textsc{Escuela Técnica Superior de Ingenierías Informática y de Telecomunicación}\\
\textsc{---}\\
Granada, Septiembre de 2019
\end{minipage}
\end{titlepage}


	% Plantilla prefacio UGR
	\thispagestyle{empty}

\begin{center}
{\large\bfseries Sistema de bajo coste para detección de dispositivos WiFi y Bluetooth}\\
\end{center}
\begin{center}
Adrián Álvarez Sáez\\
\end{center}

%\vspace{0.7cm}

\vspace{0.5cm}
\noindent{\textbf{Palabras clave}: ciudades inteligentes, sensores Bluetooth, sensores WiFi, floating car data, tráfico rodado, probe data
\vspace{0.7cm}

\noindent{\textbf{Resumen}\\

El gran avance en tecnologías móviles durante los últimos años ha propiciado que hoy en día la mayoría población utilice dispositivos móviles inalámbricos como \textit{smartphones} o \textit{smartwatchs}, estos dispositivos están emitiendo continuamente información al medio que permite identificarlos de manera única, surge la oportunidad de utilizar estos datos para analizar patrones de movimiento de la población que permitan optimizar las infraestructuras y los recursos de las ciudades con el fin de mejorar la calidad de vida de las personas y del entorno.\\

En este proyecto se ha desarrollado un sensor de bajo coste  y reducido consumo energético basado en soluciones \textit{System on chip} que permite la detección dispositivos inalámbricos en tiempo real mediante el análisis del tráfico WiFi y Bluetooth.\\

\cleardoublepage

\begin{center}
	{\large\bfseries Low cost system for detecting WiFi and Bluetooth devices}\\
\end{center}
\begin{center}
	Adrián Álvarez Sáez\\
\end{center}
\vspace{0.5cm}
\noindent{\textbf{Keywords}: smart cities, Bluetooth sensors, WiFi sensors, floating car data, road traffic, probe data
\vspace{0.7cm}

\noindent{\textbf{Abstract}\\
The advancement in mobile technologies in recent years, has led the majority of the population to use wireless mobile devices such as smartphones or smartwatchs, these devices are continuously broadcasting information that can be used to uniquely identify them, this bring us the opportunity to use these data to analyze population movement patterns that allow optimizing the infrastructure and resources of cities in order to improve the quality of life of people and the environment.\\

In this project, a low cost and low energy consumption sensor based on System on chip solutions has been developed that allows detection of wireless devices in real time by analyzing WiFi and Bluetooth traffic.\\

\cleardoublepage

\thispagestyle{empty}

\noindent\rule[-1ex]{\textwidth}{2pt}\\[4.5ex]

D. \textbf{Juan Julián Merelo Guervós}, Profesor del Departamento de
Arquitectura y Tecnología de Computadores de la Universidad de
Granada.

\vspace{0.5cm}

\textbf{Informo:}

\vspace{0.5cm}

Que el presente trabajo, titulado \textit{\textbf{Sistema de bajo coste para detección de dispositivos WiFi y Bluetooth}},
ha sido realizado bajo mi supervisión por \textbf{Adrián Álvarez Sáez}, y autorizo la defensa de dicho trabajo ante el tribunal
que corresponda.

\vspace{0.5cm}

Y para que conste, expiden y firman el presente informe en Granada a Septiembre de 2020.

\vspace{1cm}

\textbf{El/la director(a)/es: }

\vspace{5cm}

\noindent \textbf{Juan Julián Merelo Guervós}

\chapter*{Agradecimientos}


	% Índice de contenidos
	\newpage
	\tableofcontents

	% Índice de imágenes y tablas
	\newpage
	\listoffigures

	% Si hay suficientes se incluirá dicho índice
	\listoftables
	\newpage

	% Introducción
	\subfile{secciones/01_introduccion}

  % Estado del arte
	% 	1. Crítica al estado del arte
	% 	2. Propuesta
	\subfile{secciones/02_estado_del_arte}

	% Descripción del problema y hasta donde se llega
	\subfile{secciones/03_descripcion}

	\subfile{secciones/04_planificacion}

	% Análisis del problema
	% 1. Análisis de requisitos
	% 2. Análisis de las soluciones
	% 3. Solucion propuesta
	% 4. Análisis de seguridad
	\subfile{secciones/05_analisis}

	% Desarrollo bajo sprints:
	% 	1. Permitir registros y login de usuarios
	% 	2. Desarrollo del sistema de incidencias
	% 	3. Desarrollo del sistema de denuncias administrativas y accidentes
	% 	4. Desarrollo del sistema de croquis
	%   5. Instalación de la aplicación de manera automática
	\subfile{secciones/06_implementacion}

	% Presupuesto

	% Conclusiones
	\subfile{secciones/07_conclusiones}

	% Trabajos futuros



	\newpage

	\printbibliography
	%\bibliography{bibliografia}
	%\bibliographystyle{plain}

\end{document}
